% ============================================================================
% naranjas.tex  —  NARANJAS = dinero (obra propia, CC0)
% ----------------------------------------------------------------------------
% La "moneda que crece" de la serie. Plano, tierno, recoloreable, apto SVGMobject.
% Requiere: % ============================================================================
% _estilo_interes.tex  —  Estilo compartido del "mundo del capibara" (CC0)
% ----------------------------------------------------------------------------
% Paleta plana + estilo de linea + FRAME canonico para trajes/accesorios que se
% dibujan EN LAS MISMAS COORDENADAS que el capibara base
% (tikzlib/interes/capibara_kawaii.tex). Todo obra propia, CC0.
%
% Requiere: \usepackage{tikz}. NO usa texto/fuentes -> apto SVGMobject de Manim.
%
% FRAME CANONICO DE REGISTRO (\capiframe):
%   Todos los trajes/accesorios y el propio capibara se exportan con ESTE mismo
%   \useasboundingbox. Asi los SVG comparten viewBox y en Manim se pueden cargar
%   con  SVGMobject(..., should_center=False)  y quedan PERFECTAMENTE alineados
%   sobre el capibara sin tocar nada (registro automatico).
%   Con should_center=True (por defecto) el frame se ignora (Manim recorta al
%   contenido visible), asi que NO rompe usos existentes del capibara.
% ============================================================================

% --- linework del capibara (idempotente con capibara_kawaii.tex) ---
\providecolor{capicontorno}{HTML}{6E4E34}  % linea (marron calido suave)
\tikzset{
  capilinea/.style   ={draw=capicontorno, line width=2.1pt,
                       line cap=round, line join=round},
  capilineaf/.style  ={draw=capicontorno, line width=1.6pt,
                       line cap=round, line join=round}, % linea fina
}

% --- paleta de trajes/accesorios (plana, tierna) ---
\providecolor{vueltiaobase}{HTML}{F2E6C6}  % cana flecha (crema claro)
\providecolor{vueltiaosom} {HTML}{E4D2A6}  % sombra del ala
\providecolor{vueltiaopin} {HTML}{3B332B}  % "pintao" (bandas oscuras)

\providecolor{banqcorbata} {HTML}{9C3B4C}  % corbatin vino
\providecolor{banqcorbdark}{HTML}{7E2E3D}  % corbatin sombra
\providecolor{oro}         {HTML}{D2A63C}  % oro (monoculo/detalles)
\providecolor{orodark}     {HTML}{B4862A}

\providecolor{comerdelan}  {HTML}{5E86A6}  % delantal azul mercado
\providecolor{comerdeldark}{HTML}{4B6E8C}
\providecolor{comergorra}  {HTML}{CF5B45}  % gorra roja
\providecolor{comergordark}{HTML}{B24734}

\providecolor{acadboard}   {HTML}{343747}  % birrete azul-noche
\providecolor{acadbdark}   {HTML}{262838}
\providecolor{acadgafas}   {HTML}{3B332B}  % marco de gafas

\providecolor{ahorrosa}    {HTML}{F2A6BC}  % cerdito rosado
\providecolor{ahorrodark}  {HTML}{E086A0}
\providecolor{ahorrogorro} {HTML}{D9A441}  % gorro mostaza

\providecolor{naranjafill} {HTML}{F59E2E}  % naranja
\providecolor{naranjadark} {HTML}{E07C1A}  % naranja sombra
\providecolor{naranjacont} {HTML}{9A5A1E}  % contorno naranja
\providecolor{hojaverde}   {HTML}{6BA83C}  % hoja
\providecolor{hojadark}    {HTML}{4E8A2E}
\providecolor{canastamimb} {HTML}{C8925A}  % mimbre canasta
\providecolor{canastadark} {HTML}{A9743F}

% --- fondo/escena ---
\providecolor{escpared}    {HTML}{FCF4E2}  % pared crema (combina con el video)
\providecolor{escpiso}     {HTML}{EFE1C4}  % piso
\providecolor{esczocalo}   {HTML}{E4D2A6}
\providecolor{banqcol}     {HTML}{D9E3EC}  % columnas banco
\providecolor{banqcoldark} {HTML}{BFCEDC}
\providecolor{banqarco}    {HTML}{C7D6E4}
\providecolor{merctoldoA}  {HTML}{D9694E}  % toldo rojo
\providecolor{merctoldoB}  {HTML}{FBF1DC}  % toldo crema
\providecolor{mercmadera}  {HTML}{C8925A}  % madera del puesto
\providecolor{mercmaddark} {HTML}{A9743F}
\providecolor{fincacielo}  {HTML}{CDE7F0}  % cielo
\providecolor{fincapasto}  {HTML}{9FCB6B}  % pasto
\providecolor{fincapastod} {HTML}{7FB050}
\providecolor{fincasol}    {HTML}{F6C64B}  % sol
\providecolor{fincamata}   {HTML}{6BA83C}

% FRAME CANONICO (capibara + trajes) — simetrico en x; contiene sombreros altos
% (hasta y~2.05) y la alcancia lateral (hasta x~2.65).
\newcommand{\capiframe}{\useasboundingbox (-2.75,-2.10) rectangle (2.75,2.25);}

%
% Uso:
%   \naranjaunit           % una naranja (unidad monetaria), centrada en (0,0), r~0.5
%   \naranjaunit[esc]      % --- ver macro \naranjaen para colocar/escalar ---
%   \naranjaen{x}{y}{s}    % una naranja en (x,y) escalada s (para armar pilas)
%   \naranjapila           % PILA/monton (piramide de 10) -> "muchas"
%   \naranjacanasta        % canasta/saco lleno de naranjas
%
% NARANJA UNIDAD: circulo naranja + brillo + ombligo (estrellita) + hoja.
% Diametro ~1.0 (radio 0.5). Recolorear: \colorlet{naranjafill}{...}.
% ============================================================================

% --- una naranja centrada en (cx,cy), escala s (1 = radio 0.5) ---
\newcommand{\naranjaen}[3]{%
  \begin{scope}[shift={(#1,#2)}, scale=#3]
    \path[capilineaf, fill=naranjafill] (0,0) circle[radius=0.50];
    % gajo/sombra inferior
    \path[fill=naranjadark] (0,-0.06) circle[radius=0.44];
    \path[fill=naranjafill] (0,0.02) circle[radius=0.44];
    % ombligo (estrellita) arriba, hacia la hoja
    \path[fill=naranjadark] (0.00,0.34) circle[radius=0.05];
    % brillo
    \path[fill=vueltiaobase] (-0.20,0.18) ellipse[x radius=0.12, y radius=0.16];
    % hoja
    \path[capilineaf, fill=hojaverde]
        (0.06,0.44) .. controls (0.30,0.70) and (0.02,0.72) .. (-0.06,0.50)
        .. controls (-0.02,0.58) and (0.10,0.56) .. (0.06,0.44) -- cycle;
    \draw[capilineaf, draw=hojadark] (0.02,0.48) -- (0.14,0.62);
  \end{scope}
}

% --- una naranja unidad centrada (comodin) ---
\newcommand{\naranjaunit}{\naranjaen{0}{0}{1}}

% --- PILA / MONTON: piramide de 10 naranjas (4-3-2-1) = "muchas" ---
%   Base ancha, crece hacia arriba: ideal para "1 -> varias -> muchas".
%   Centrada en x=0; base en y~-0.55. Dibuja de atras hacia adelante.
\newcommand{\naranjapila}{%
  \begin{scope}
    % fila 4 (base)
    \naranjaen{-1.35}{-0.55}{0.9}\naranjaen{-0.45}{-0.55}{0.9}%
    \naranjaen{ 0.45}{-0.55}{0.9}\naranjaen{ 1.35}{-0.55}{0.9}%
    % fila 3
    \naranjaen{-0.90}{ 0.20}{0.9}\naranjaen{ 0.00}{ 0.20}{0.9}%
    \naranjaen{ 0.90}{ 0.20}{0.9}%
    % fila 2
    \naranjaen{-0.45}{ 0.95}{0.9}\naranjaen{ 0.45}{ 0.95}{0.9}%
    % fila 1 (cima)
    \naranjaen{ 0.00}{ 1.70}{0.9}%
  \end{scope}
}

% --- CANASTA / SACO de naranjas (montoncito rebosante en canasta de mimbre) ---
\newcommand{\naranjacanasta}{%
  \begin{scope}
    % naranjas asomando por atras del borde
    \naranjaen{-0.62}{0.30}{0.85}\naranjaen{0.00}{0.52}{0.9}%
    \naranjaen{ 0.66}{0.30}{0.85}\naranjaen{-0.20}{0.66}{0.6}%
    \naranjaen{ 0.30}{0.66}{0.6}%
    % cuerpo de la canasta (trapecio redondeado)
    \path[capilinea, fill=canastamimb]
        (-1.15,0.30) -- (1.15,0.30) -- (0.92,-1.02) -- (-0.92,-1.02) -- cycle;
    % borde superior
    \path[capilinea, fill=canastadark] (0,0.30) ellipse[x radius=1.15, y radius=0.22];
    \path[fill=canastamimb] (0,0.34) ellipse[x radius=1.02, y radius=0.15];
    % tejido (lineas de mimbre)
    \draw[capilineaf, draw=canastadark] (-0.98,-0.10) -- (0.98,-0.10);
    \draw[capilineaf, draw=canastadark] (-0.90,-0.52) -- (0.90,-0.52);
    \foreach \x in {-0.7,-0.35,0,0.35,0.7}
      \draw[capilineaf, draw=canastadark] (\x,0.28) -- (\x,-1.00);
  \end{scope}
}
